\origpage{78}{N. Robertson, P. Seymour / Journal of Combinatorial Theory, Series B 89 (2003) 77--108}
\noindent conjecture, first to prove that the class of finite hypergraphs with edges of size 2 or 3 drawable on a fixed surface is well-quasi-ordered, and secondly to derive from this that the class of all finite graphs is well-quasi-ordered. It will also be needed in later papers, again in a hypergraph form, to prove Nash-Williams' ``immersion'' conjecture [1]. We shall therefore formulate it completely in terms of hypergraphs.

The paper is organized as follows. Section 2 contains basic definitions and results about tree-decompositions and tangles. Sections 3 and 4 develop the relation between the kinds of decomposition relative to a tangle that we need. In Section 5 we introduce patchworks, which enable us to define minors of hypergraphs, and develop some lemmas about them. The main result is stated and proved in Section 6, and Section 7 contains a lemma which is often useful in applying the theorem.

Thus this work falls into two parts. Sections 2--4 are about how to convert information about the local structure of a hypergraph relative to each of its high-order tangles, into a linked tree-decomposition whose pieces (the nodes of the tree) correspond to the high-order tangles, and still have the same local structure (more or less---we may have to grow the pieces to make them fit together by adding on subhypergraphs of bounded tree-width). Here the tree-decompositions use unrooted trees; there is no reason to fix a root for the tree, and if we did so the results would appear most unnatural. The second half, Sections 5 and 6, mostly concerns well-quasi-ordering, and in that topic we have to use rooted trees; we have to do complicated inductions concerning the sizes of these trees, and it is very important to fix a root of the tree. When we do so, for each piece of the tree-decomposition, there is not symmetry between its neighbouring pieces any more; one is in the direction of the root, and has to be treated differently. When we lop off the arms of the tree-decomposition growing out from a given piece, and replace these arms by new hyperedges marking where the arms used to attach (which is what we mean by the local structure at the node of the tree), it is convenient to lop off the ``root arm'' in a different way; instead of replacing it by a new hyperedge, we simply label the vertices where it used to attach and call them roots of the hypergraph. And also, when we lop off the ``non-root'' arms, we need to remember not only the set of vertices where the arm used to attach, but also which of these vertices was which; we need to remember an ordered set. So the new hyperedge replacing the arm will have to be equipped with a linear order of its vertex set. The point is that half-way through the paper, suddenly our trees become rooted trees or ``arborescences'', and the hypergraphs develop roots, and their hyperedges become ordered sets of vertices. This is most confusing when it happens (particular since we have to redefine all our terms for rooted trees and rooted hypergraphs, and there is not quite an exact correspondence), and we hope it will help the reader to be warned ahead of time.

\sectionhead{2. Hypergraphs, tangles and tree-decompositions}

For the purposes of this paper, a \textit{hypergraph} $G$ consists of a finite set $V(G)$ of \textit{vertices}, a finite set $E(G)$ of \textit{edges}, and an incidence relation between them. The
